\chapter{Phase 4: nonlinear equations}\label{ch:cII4}
\chapterbrief
{Part~I, Chapter~\ref{ch:cI4} (A25--A31), and its notation: Burgers flows $S$, $S_B$, $S_A$, the Lagrangian map $X$, $J$, the Hopf--Lax functional $\Phi_x$, and the dictionary setting.}
{to prove the tools behind A26--A28 (comparison principle, pre-shock bounds), the blow-up at the shock time, the jumps after the shock, the $\sqrt\nu$ splitting error, the KdV--Burgers remainder, and the remaining identification results (the negative theorem for complete dictionaries, closed forms, the degree bound for Burgers).}
{Dim.~B61--B72, followed by the list of nonlinear results whose long proofs are not reproduced here.}

%=====================================================================
\section{Tools before the shock}

\begin{dimcard}{B61}{S4T-L0}{Comparison principle}{\proved}{}
Let $b,c$ be continuous and real on $[0,T]\times\mathbb T$, $|c|\le K$, $\nu\ge0$, and $Lg=g_t+bg_x+cg-\nu g_{xx}$. (a) If $Lg\ge0$ and $g(0)\ge0$ then $g\ge0$. (b) If $|Lg|\le F(t)$ then
\[
\|g(t)\|_\infty\le e^{Kt}\Big(\|g(0)\|_\infty+\int_0^tF\Big).
\]
\end{dimcard}
\begin{proof}
\step{1}{(a) First touching point}
For $\theta>0$ put $z=g+\theta e^{(K+1)t}$, so $Lz>0$ and $z(0)>0$. At a first point $(t_1,x_1)$ with $z=0$: $z_t\le0$, $z_x=0$, $z_{xx}\ge0$, so $Lz\le0$, a contradiction. Let $\theta\to0$.
\step{2}{(b)}
$Y(t)=e^{Kt}(\|g(0)\|+\int_0^tF)$ satisfies $Y'\ge KY+F$; apply (a) to $Y\mp g$.
\whyworks{a parabolic operator cannot create a new minimum.}
\end{proof}
\uses{none.}

\begin{dimcard}{B62}{S4T-L1}{Derivative bounds before the shock}{\proved}{}
For real nonconstant $u_0$, $M_0=\max u_0'$, $0<T<t^*$, and $w$ the viscous or inviscid Burgers solution: $-1/(t^*-t)\le\partial_xw\le M_0$, $\|w\|_\infty\le\|u_0\|_\infty$, and $\|\partial_x^jw(t)\|_\infty\le Y_j(t)$ with $Y_j$ independent of $\nu$.
\end{dimcard}
\begin{proof}
\step{1}{The slope}
$p=\partial_xw$ solves $p_t+wp_x+p^2-\nu p_{xx}=0$. Apply B61(a) to $M_0-p$ (source $p^2\ge0$) and to $p-\vartheta$ with $\vartheta=-1/(t^*-t)$, which solves $\vartheta'=-\vartheta^2$.
\step{2}{Higher derivatives}
$p_j=\partial^jw$ solves $p_{j,t}+wp_{j,x}+(j+1)p_1p_j-\nu p_{j,xx}=-\sum_{i=2}^{j-1}\binom jip_ip_{j+1-i}$; B61(b) and induction give $\nu$-independent bounds.
\whyworks{before the shock every derivative obeys a transport equation with bounded coefficients.}
\end{proof}
\uses{B61.}

\begin{dimcard}{B63}{S4T-P3}{The closed form blows up at the shock time}{\pw}{}
If $u_0'$ has a unique minimum $a^*$ with $\kappa=u_0'''(a^*)>0$ and $\eta=1-t/t^*$, the first-order coefficients of A26 satisfy, as $t\uparrow t^*$,
\[
\eta^{3/2}\sup_x|E_1^{\mathrm{add}}(t)|\to\frac{3\sqrt3}{16}\sqrt{2t^*\kappa}.
\]
\end{dimcard}
\begin{proof}
\step{1}{Local form}
With $y=a-a^*$: $u_0''(a)=\kappa y(1+O(y))$ and $J=\eta+\tfrac{t\kappa}{2}y^2(1+O(y))$.
\step{2}{Scaling}
With $y=z\sqrt{\eta/c_t}$, $c_t=t\kappa/2$, the quantity $|u_0''|/J^2$ equals $\frac{\kappa}{\sqrt{c_t}}\eta^{-3/2}\,z/(1+z^2)^2$ up to $1+O(\eta^{1/4})$; the maximum of $z/(1+z^2)^2$ is $3\sqrt3/16$. Away from the window the function is smaller.
\whyworks{the Jacobian vanishes quadratically at the cusp, and the coefficient is a ratio of a linear to a squared quadratic quantity.}
\end{proof}
\uses{A26.}

%=====================================================================
\section{After the shock}

\begin{dimcard}{B64}{S4T-T3}{After the shock: jumps and the splitting order}{\modulo{[F1]}}{}
For smooth nonconstant $u_0$, $t>t^*$ and every $\nu>0$: (i) $\Eint^{\mathrm{add}}(t)$ has a jump $+D$ at each shock, independent of $\nu$, so $\|\Eint^{\mathrm{add}}\|_\infty\ge D/2$ and $\Eint^{\mathrm{add}}\notin H^{1/2}$; (ii) the same for the splitting error $B$ then $A$; (iii) the splitting error $A$ then $B$ is smooth but $\liminf_\nu\nu^{-1/2}\|E^{\mathrm{mult}}_{AB}\|_{L^1}\ge2D\sqrt{t/\pi}$.
\end{dimcard}
\begin{proof}
\step{1}{(i)}
$\Eint^{\mathrm{add}}=S(t)u_0-S_B(t)u_0-S_A(t)u_0+u_0$: the viscous, heat and initial terms are smooth, while $S_B(t)u_0$ jumps by $-D$ at the shock (Lax--Oleinik).
\step{2}{(iii)}
$E^{\mathrm{mult}}_{AB}=S(t)u_0-S_A(t)S_B(t)u_0$; heat-smoothing a jump of size $D$ during time $t$ leaves an error of $L^1$ size $\ge2D\sqrt{\nu t/\pi}$, and the viscous solution is $O(\nu|\log\nu|)$-close to the shock profile.
\whyworks{the inviscid term alone produces a discontinuity that no amount of viscosity in the other term can remove.}
\end{proof}
\uses{[F1], [F2].}

\begin{dimcard}{B65}{OPF5D-T6}{The vanishing-viscosity $L^1$ bound}{\proved{} (mod [F1] for the corollary)}{prop:splitting}
For every $t,\nu>0$,
\[
\|S(t)u_0-v(t)\|_{L^1(\mathbb T)}\le8\pi\sqrt{\frac{2\nu\log(1+t\|u_0'\|_\infty)}{t}}.
\]
Consequently $\|E^{\mathrm{mult}}_{AB}(t)\|_{L^1}=\Theta(\sqrt\nu)$ for $t>t^*$, and $\|\Eint^{\mathrm{add}}(t)\|_{L^1}=O(\sqrt\nu)$.
\end{dimcard}
\begin{proof}
\step{1}{Convex potentials}
$G_\nu(x)=2\nu\log\int\exp\big((xy/t-h(y))/2\nu\big)\dd y$ and $G_0(x)=\sup_y(xy/t-h(y))$, $h=U_0+y^2/2t$, are convex, and $u-v=G_0'-G_\nu'$.
\step{2}{Laplace gap}
$G_\nu-G_0=2\nu\log I(x)$ with $I\ge c_0=\sqrt{4\pi\nu/M}$, $M=\|u_0'\|_\infty+1/t$; put $\varepsilon=G_\nu-c-G_0\ge0$.
\step{3}{Strong convexity and averaging}
$G_0$ is $(1/t^2M)$-strongly convex, so each $D_x(y)$ is quadratic-from-below in $x$; Fubini over a period and Jensen give $\int_{\mathbb T}\varepsilon\le4\pi\nu\log(tM)$.
\step{4}{Derivatives of close convex functions}
With $f=G_\nu-c$, $g=G_0$, $f-g=\varepsilon\ge0$ and $h>0$, convexity gives
\[
|f'-g'|(x)\le g'(x+h)-g'(x-h)+\frac{\varepsilon(x+h)+\varepsilon(x-h)}{h}.
\]
Over one period $\int(g'(x+h)-g'(x-h))\dd x=4\pi h/t$, so
\[
\|u-v\|_{L^1}\le\frac{4\pi h}{t}+\frac{8\pi\nu\log(tM)}{h},
\]
and $h=\sqrt{2\nu t\log(tM)}$ gives the bound.
\step{5}{The splitting error}
$E^{\mathrm{mult}}_{AB}=(u-v)-(S_Av-v)$ with $\|S_Av-v\|_{L^1}\le2\sqrt{\nu t/\pi}\,\mathrm{TV}(v(t))$; the lower bound is B64(iii).
\whyworks{the viscous and inviscid solutions are derivatives of two convex potentials that differ by a Laplace correction of size $\nu\log$.}
\end{proof}
\uses{[F1], [F2].}

\begin{dimcard}{B66}{OPR2-T5}{KdV--Burgers before the shock}{\modulo{local well-posedness}}{}
For $u_t+uu_x=\nu u_{xx}-du_{xxx}$, $k\ge2$, $T<t^*$, and small $\nu+|d|$:
\[
\sup_{t\le T}\big\|u(t)-v(t)-\nu\delta(t)-d\,\delta_d(t)\big\|_{H^k}\le C_k(\nu+|d|)^2,
\]
so $\Eint^{\mathrm{add}}$ has an explicit first-order form in $(\nu,d)$.
\end{dimcard}
\begin{proof}
\step{1}{Energy estimate for $e=u-v$}
$e_t+(ve)_x+ee_x-\nu e_{xx}+de_{xxx}=\nu v_{xx}-dv_{xxx}$. In $H^k$ the dispersive term is antisymmetric, the viscous term dissipates, and the transport terms are bounded by $C\|e\|_{H^k}^2$.
\step{2}{Bootstrap}
$(\sqrt Y)'\le C\sqrt Y+G(\nu+|d|)$ keeps $\|e\|_{H^k}\le C(\nu+|d|)$ on $[0,T]$; the same for $e-\nu\delta-d\delta_d$ gives the square.
\whyworks{before the shock both small terms are regular perturbations of the inviscid flow.}
\end{proof}
\uses{B62.}

\begin{dimcard}{B67}{S4T-P5}{Fisher--KPP: the jet}{\proved}{}
For $u_t=\nu u_{xx}+ru(1-u)$ with terms $A=\nu u_{xx}$, $B=ru(1-u)$:
\[
[s^2]\,\hat E^{\mathrm{add}}=\tfrac12\big(A'[B]+B'[A]\big)=r\nu\big(u_0''-u_0'^2-2u_0u_0''\big);
\]
the estimate of the remainder is not uniform in time.
\end{dimcard}
\begin{proof}
\step{1}{The two derivatives}
$A'[w]=\nu w_{xx}$, $B'[w]=r(1-2u_0)w$. So $A'[B]=r\nu(u_0-u_0^2)''=r\nu(u_0''-2u_0'^2-2u_0u_0'')$ and $B'[A]=r\nu(1-2u_0)u_0''$.
\step{2}{Half the sum}
$\tfrac12\big(A'[B]+B'[A]\big)=r\nu\big(u_0''-u_0'^2-2u_0u_0''\big)$.
\whyworks{A25 at order two.}
\end{proof}
\uses{A25.}

\begin{example}{The reference datum $u_0=\varepsilon\sin x$ \led{OPA-C1}}
Here $t^*=1/\varepsilon$, there is one shock at $x=\pi$ for $t>t^*$, symmetric ($\rho=0$), and A27 gives
\[
\Lambda(t)=2\log\frac{1}{\varepsilon t-1}+4\log2+o(1)\qquad(t\downarrow t^*),
\]
consistent with the uniform rate of A28.
\end{example}

%=====================================================================
\section{Identification in dictionaries}

\begin{dimcard}{B68}{S4ID-P1}{The literal statement holds only in the minimal dictionary}{\proved}{}
For $F=y_t-\sigma y_n-N'$ with $\sigma\ne0$ and $\deg N=2$ (Burgers, Fisher--KPP, KdV, KdV--Burgers): (a) $F\in V_{d,k}$ iff $d\ge2$ and $k\ge n$; (b) if $(d,k)\ne(2,n)$ then $\dim([F]\cap V_{d,k})\ge2$; (c) $[F]\cap V_{2,n}=\C F$.
\end{dimcard}
\begin{proof}
\step{1}{(b) Two independent consequences}
If $k\ge n+1$: $\partial_xF\in V_{d,k}$ contains $y_{n+1}$, which $F$ does not. If $d\ge3$: $yF$ contains $y\,y_t$, which $F$ does not.
\step{2}{(c) Elimination}
In $V_{2,n}$, the generators involving time derivatives other than $y_t$ have normal forms of distinct orders $na+b\ge n+1$ with nonzero leading coefficients $\sigma^a$, so they must cancel; what remains is a multiple of $F$ by degree count.
\whyworks{a complete dictionary always contains derivatives or multiples of the equation.}
\end{proof}
\uses{A29.}

\begin{dimcard}{B69}{S4ID-T4}{Closed forms do not identify; identifying solutions exist}{\pw{} ((iii) mod H\"older, Bona--Smith \cf)}{}
(i) If $u$ lies in a finite tower of field extensions of $\C(t,x)$ (closed forms), then $u,u_x,\dots,\partial_x^ru$ are algebraically dependent, so $\mathcal J_u\ne0$ and $u$ does not identify its equation in any dictionary containing that relation. (ii) For $u=P/Q$ in $t,x$ and exponentials, an explicit degree $D^*=(n+1)c^n-n$ of such a relation exists. (iii) For a datum $g$ that satisfies no algebraic ODE (for example $\Gamma$ on an interval), the Burgers or KdV solution has $\mathcal J_u=0$ and identifies its equation modulo consequences in every dictionary.
\end{dimcard}
\begin{proof}
\step{1}{(i) Transcendence degree}
Each step of the tower adds at most a bounded transcendence degree, so the field generated by $u$ and its $x$-derivatives has finite transcendence degree $r$; $r+1$ derivatives are dependent. A29(c) concludes.
\step{2}{(iii) A transcendental datum}
At $t=0$ any relation in the $x$-jet would be an algebraic ODE for $g$; by continuity a relation for $u$ restricts to one for $g$. H\"older's theorem says $\Gamma$ satisfies none \cf.
\whyworks{closed forms live on finite-dimensional varieties of jets; differentially transcendental data do not.}
\end{proof}
\uses{A29.}

\begin{dimcard}{B70}{OP-T7}{Burgers $N$-shocks: no relation of degree $\le N$}{\proved}{thm:burgersN}
\phantomsection\label{dim:S4ID-K1}%
Under the hypotheses of A30, no nonzero polynomial of degree $\le N$ in the jet of $u$ (of any order) vanishes on $u$. Hence, for every order $K\ge N$, $u$ identifies Burgers in the evolutionary dictionary iff $d\le N$.
\end{dimcard}
\begin{proof}
\step{1}{Cumulants}
By A30, Steps~1--2, the jet of $w$ is the sequence of cumulants $\kappa_n$ of the moments $r_j$, polynomial in $r_1,\dots,r_N$, which are algebraically independent.
\step{2}{Weights}
Give $r_j$ weight $j$. The top-weight part of $\kappa_n$ is a truncated cumulant; a relation $P(\kappa)=0$ of degree $\le N$ has a nonzero top-weight part that must vanish identically.
\step{3}{Confluent limit}
The top-weight parts are cumulants of a measure with $N$ atoms; in the confluent limit they become power sums of $N$ variables, and power sums $p_1,\dots,p_N$ are algebraically independent (fundamental theorem of symmetric polynomials), which excludes any relation of degree $\le N$.
\whyworks{$N$ atoms leave $N$ independent power sums, so the first relation needs degree $N+1$.}
\end{proof}
\uses{A30.}

\begin{dimcard}{B71}{F4L-P3}{Analyticity radius of Burgers in the amplitude}{\pw}{}
For Burgers with datum $\varepsilon U_0'$, the solution is analytic in $\varepsilon$ for $|\mathrm{Im}\,\varepsilon|<2\pi\nu/\mathrm{osc}(U_0)$, uniformly in time, and this is optimal. \textbf{(KIEFFER-C)} The ``infinite radius'' claimed in a cited preprint is not proved there: its own formula has denominator zeros at finite $\varepsilon$.
\end{dimcard}
\begin{proof}
\step{1}{A moment generating function}
By Cole--Hopf, $\varphi_\varepsilon(t,x)=\mathbb E\,e^{sX}$ with $X=U_0(Y)$, $Y$ distributed by the heat kernel, $s=-\varepsilon/2\nu$; $X$ takes values in an interval of length $w=\mathrm{osc}\,U_0$.
\step{2}{The strip lemma}
For a random variable with values in an interval of length $w$, $M(s)=\mathbb E e^{sX}$ has no zeros for $|\mathrm{Im}\,s|<\pi/w$ (L-S), and $u=-2\nu\partial_x\log\varphi$ is analytic there. The symmetric Bernoulli law attains the bound (the $\tanh$ case).
\whyworks{Cole--Hopf turns Burgers into the logarithm of a moment generating function, whose zeros are kept away by the bounded range.}
\end{proof}
\uses{[F2].}

\begin{dimcard}{B72}{S4ID-L4}{Burgers in Hammerstein form}{\pw}{}
With $\ring$ acting on differential polynomials, and $F_B=(\tau-\nu\xi^2)y+\tfrac\xi2y^2$: $[F_B]\cap(\ring y\oplus\ring y^2)=\ring F_B$.
\end{dimcard}
\begin{proof}
\step{1}{Degree splitting}
Elements of $[F_B]$ are $\sum_jL_j\,\partial^jF_B$ with differential-polynomial coefficients; those of degree $\le2$ in $y$ with linear--quadratic structure must have constant-coefficient multipliers, because any $y$-dependent coefficient raises the degree.
\whyworks{in the Hammerstein form the linear and quadratic parts cannot be mixed without leaving the space.}
\end{proof}
\uses{A29.}

%=====================================================================
\section{Nonlinear results whose proofs are not reproduced}

The following results are proved in the research files named in the last column. Their proofs rely on long computations (scattering theory, Hirota bilinear forms, computer-assisted syzygies). They are listed with their labels so that the catalogue is complete as a list; the proofs are not reproduced in this edition.

\listcard{INT-T1}{KdV in the class $\mathcal H_M$: identification by the scattering data; (A) and unique continuation \proved; (B), (C), INT-T2 \condl{} / \proved{} at fixed time (OPF5D).}{see label}{\texttt{work/S4-INTEG} \S1.1, \texttt{work/OP-F5D}.}
\listcard{INT-P5, INT-P6}{Hirota $N$-solitons: the interaction is pairwise, integrates to $0$, costs $O(N^3)$.}{\proved}{\texttt{work/S4-INTEG}.}
\listcard{OP-P17, OP-P18}{Lax and zero-curvature relations for all $n$ (sign at $n=0$ corrected).}{\proved}{\texttt{work/OP-F5D} (OPF5D-T1, T2).}
\listcard{OPF5D-R2}{NLS identification by scattering data without multiplicity.}{\falsest{} mod [DP]}{\texttt{work/OP-F5D}.}
\listcard{OPF5D-T7}{Lax--Oleinik equals Kru\v zkov entropy solution ([F1]).}{\modulo{Kru\v zkov}}{\texttt{work/OP-F5D}.}
\listcard{OPF5A-T3, T4, C1}{OPA-T1 with multiple points; genericity of $(H_t)$; global uniform rate for generic data (several births \condl).}{\proved{} mod [F1]}{\texttt{work/OP-F5A}.}
\listcard{OPF5C-T1, L1--L8, P7, P9, T6}{Structure of Novikov's variety; upper bounds for all $N$; nine exact syzygy degrees $\delta_N(K)$ (computer-assisted).}{\proved{} mod (Nov$'$)}{\texttt{work/OP-F5C}; the main line is in A31.}
\listcard{S4ID-L2, L5, C1, C2, P5--P7}{Normal forms and elimination lemmas; evolutionary versus complete dictionaries.}{\proved}{\texttt{work/S4-ID}.}
\listcard{OP-C11, OP-T1, OP-E3, OP-C3$'$, OP-L2..P14}{Complete criteria for the Burgers family; KdV examples and obstructions.}{\proved}{\texttt{work/OP-NL}.}
\listcard{ALG-T1, MODREL-T1, F4LIN-T2}{Identification in the full ring for analytic $u$; moment relations; tropical interaction $4\nu\ln2$.}{\proved}{\texttt{work/S4-INTEG}.}
\listcard{OPR2-P4, P6, O8, O9$'$}{Novikov relation, obstructions to INT-Q1.}{\proved{} / mod Novikov}{\texttt{work/OP-R2}.}
\listcard{S4ID-P8, OPR2-P7, OPF5D-P14, P15}{Nonlinear asymptotics with derivatives (Deift--Zhou).}{\condl}{\texttt{work/S4-ID}, \texttt{work/OP-F5D}.}
\listcard{INT-Q1, OPF5C-Q1}{Analytic KdV data with reflection; general $\delta_N(K)$.}{\openq}{}
